% ====================================================================
%  徐森林、金亚东、薛春华《数学分析》（第三册）
%  \S16.4  Fourier 级数的 Cesaro 求和
%  源：原书印刷页 299–304（PDF 物理页 305–310）
%  约定：全西文标点；行间公式用一对方括号；多行公式用 align 或 aligned；
%        跨册引用一律用 extref；数学模式内的中文必须包 text
% ====================================================================
\section{Fourier 级数的 Cesaro 求和}\label{sec:16.4}

Dini 的收敛判别法指出, 除了要求 $f(x)$ 在 $x_{0}$ 处连续外, 还要求 $f(x)$ 在 $x_{0}$ 处有一阶导数, 或有广义的左右导数. 自然要问: 仅有 $f(x)$ 的连续性, 是否能保证它的 Fourier 级数收敛于自己? $1876$ 年, Du Bois-Reymond 举出了一个连续函数, 它的 Fourier 级数在若干点是发散的. 从而否定地回答了上述提出的问题.

转而, 人们不再在连续函数上加条件, 而是去改进级数收敛的定义, 使得在新的收敛意义下, 连续函数的 Fourier 级数能收敛于自己.

设 $\displaystyle\sum_{n=1}^{\infty}a_{n}$ 为一个无穷级数, $S_{n}=\displaystyle\sum_{k=1}^{n}a_{k}$ 为它的部分和. 我们曾定义, 如果 $\lim\limits_{n\to+\infty}S_{n}=S$, 就称级数 $\displaystyle\sum_{n=1}^{\infty}a_{n}$ 收敛于 $S$ 且 $S$ 为它的和. 这个定义既自然又直观, 它与人们通常的认识是一致的. 但它的不足之处是一些很简单的级数, 在上述意义下却没有和. 例如级数
\[
  \sum_{n=1}^{\infty}(-1)^{n-1}=1-1+1-1+\cdots
\]
就是如此.

新给出的收敛定义, 必须使得在原来意义下收敛的级数, 在新的意义下仍然收敛, 而且有相同的和; 而一些在原来意义下发散的级数, 在新的意义下却是收敛的. 换句话说, 新的定义必须比原来的定义能使更多的级数有和. 下面介绍的 Cesaro 求和就能满足这样的要求.

\begin{definition}\label{def:16.4.1}
设 $\displaystyle\sum_{n=1}^{\infty}a_{n}$ 为一个无穷级数, $\{S_{n}\}$ 为它的部分和数列. 如果 $\{S_{n}\}$ 的算术平均值数列
\[
  \sigma_{n}=\frac{S_{1}+S_{2}+\cdots+S_{n}}{n},\quad n=1,2,\cdots
\]
是一个收敛数列, 即 $\lim\limits_{n\to+\infty}\sigma_{n}=\sigma$, 则称级数 $\displaystyle\sum_{n=1}^{\infty}a_{n}$ 在 \textbf{Cesaro 意义下收敛}或\textbf{均值意义下收敛}, $\sigma$ 就称为级数 $\displaystyle\sum_{n=1}^{\infty}a_{n}$ 的 \textbf{Cesaro 和}, 记为 $\displaystyle\sum_{n=1}^{\infty}a_{n}=\sigma\ (\mathrm{C})$, 此时称级数\textbf{可以 Cesaro 求和}.
\end{definition}

\begin{theorem}\label{thm:16.4.1}
如果级数 $\displaystyle\sum_{n=1}^{\infty}a_{n}$ 在通常意义下收敛于 $S$, 即 $\lim\limits_{n\to+\infty}S_{n}=S$, 则 $\{S_{n}\}$ 的算术平均值数列
\[
  \sigma_{n}=\frac{S_{1}+S_{2}+\cdots+S_{n}}{n}
\]
也收敛于 $S$, 即 $\displaystyle\sum_{n=1}^{\infty}a_{n}$ 在 Cesaro 意义下也收敛, 而且有相同的和 $S$. 但反之不成立, 即存在级数在 Cesaro 意义下收敛, 而在通常意义下是发散的.
\end{theorem}

\begin{proof}
根据\extref{例 1.1.15} 知, 从 $\lim\limits_{n\to+\infty}S_{n}=S$ 立即推得
\[
  \lim_{n\to+\infty}\sigma_{n}=\lim_{n\to+\infty}\frac{S_{1}+S_{2}+\cdots+S_{n}}{n}=S.
\]
但反之并不一定成立. 例如, $\displaystyle\sum_{n=1}^{\infty}(-1)^{n-1}$, $a_{n}=(-1)^{n-1}$,
\[
  S_{n}=\begin{cases}
    1,&n\text{ 为奇数},\\
    0,&n\text{ 为偶数},
\end{cases}
\]
在通常意义下 $\{S_{n}\}$ 是发散的. 但由
\[
  \sigma_{2k}=\frac{S_{1}+S_{2}+\cdots+S_{2k}}{2k}=\frac{k}{2k}=\frac{1}{2},
\]
\[
  \sigma_{2k+1}=\frac{S_{1}+S_{2}+\cdots+S_{2k-1}}{2k+1}=\frac{k+1}{2k+1}\to\frac{1}{2}\quad(k\to+\infty),
\]
得
\[
  \sum_{n=1}^{\infty}(-1)^{n-1}=\lim_{n\to+\infty}\sigma_{n}(\mathrm{C})=\frac{1}{2}(\mathrm{C}).
\]
\end{proof}

现在将 Cesaro 求和法用到 Fourier 级数上去. 可以证明: 连续函数 $f(x)$ 的 Fourier 级数在 Cesaro 意义下一定收敛于自己.

设 $f(x)$ 是周期为 $2\pi$ 且在 $[-\pi,\pi]$ 上可积或绝对可积的函数.
\[
  f(x)\sim\frac{a_{0}}{2}+\sum_{n=1}^{\infty}\left(a_{n}\cos nx+b_{n}\sin nx\right),
\]
它的部分和为
\[
  S_{n}(x_{0})=\frac{1}{\pi}\int_{0}^{\pi}\left[f(x_{0}+t)+f(x_{0}-t)\right]\frac{\sin\left(n+\frac{1}{2}\right)t}{2\sin\frac{t}{2}}\mathrm{d}t.
\]
利用三角恒等式
\[
  \sum_{k=0}^{n-1}\sin\left(k+\frac{1}{2}\right)t=\frac{\sin^{2}\frac{nt}{2}}{\sin\frac{t}{2}}
\]
得到 $S_{0}(x_{0})$, $S_{1}(x_{0})$, $\cdots$, $S_{n-1}(x_{0})$ 的算术平均
\begin{align*}
  \sigma_{n}(x_{0})&=\frac{1}{n}\sum_{k=0}^{n-1}S_{k}(x_{0})\\
  &=\frac{1}{n\pi}\int_{0}^{\pi}\left[f(x_{0}+t)+f(x_{0}-t)\right]\sum_{k=0}^{n-1}\frac{\sin\left(k+\frac{1}{2}\right)t}{2\sin\frac{t}{2}}\mathrm{d}t\\
  &=\frac{1}{2n\pi}\int_{0}^{\pi}\left[f(x_{0}+t)+f(x_{0}-t)\right]\left[\frac{\sin\frac{nt}{2}}{\sin\frac{t}{2}}\right]^{2}\mathrm{d}t.
\end{align*}
特别地, 取 $f(x)=1$, 则 $S_{n}(x)=1$, $\sigma_{n}(x_{0})=1$, 代入上式得到
\[
  \frac{1}{n\pi}\int_{0}^{\pi}\left(\frac{\sin\frac{nt}{2}}{\sin\frac{t}{2}}\right)^{2}\mathrm{d}t=1.
\]

\begin{theorem}[Fej\'er 定理]\label{thm:16.4.2}
设 $f(x)$ 是周期为 $2\pi$ 且在 $[-\pi,\pi]$ 上可积或绝对可积的函数. 如果 $f(x)$ 在 $x_{0}$ 处有左、右极限 $f(x_{0}-0)$ 与 $f(x_{0}+0)$, 则它的 Fourier 级数在 $x_{0}$ 处的 Cesaro 和为
\[
  \frac{1}{2}\left[f(x_{0}+0)+f(x_{0}-0)\right].
\]
特别地, 当 $f(x)$ 在 $x_{0}$ 处连续时, 它的 Fourier 级数的 Cesaro 和为 $f(x_{0})$.
\end{theorem}

\begin{proof}
记 $s=\dfrac{1}{2}[f(x_{0}+0)+f(x_{0}-0)]$, $\varphi(t)=f(x_{0}+t)+f(x_{0}-t)-2s$, 则
\begin{align*}
  \sigma_{n}(x_{0})-s&=\frac{1}{2n\pi}\int_{0}^{\pi}\left[f(x_{0}+t)+f(x_{0}-t)-2s\right]\left[\frac{\sin\frac{nt}{2}}{\sin\frac{t}{2}}\right]^{2}\mathrm{d}t\\
  &=\frac{1}{2n\pi}\int_{0}^{\pi}\varphi(t)\left[\frac{\sin\frac{nt}{2}}{\sin\frac{t}{2}}\right]^{2}\mathrm{d}t.
\end{align*}
由于 $f(x)$ 的左、右极限 $f(x_{0}-0)$ 与 $f(x_{0}+0)$ 都存在, 故 $\forall\varepsilon>0$, $\exists\delta\in(0,\pi)$, 当 $0<t<\delta$ 时,
\[
  |f(x_{0}+t)-f(x_{0}+0)|<\frac{\varepsilon}{2},\quad|f(x_{0}-t)-f(x_{0}-0)|<\frac{\varepsilon}{2},
\]
\begin{align*}
  |\varphi(t)|&=|f(x_{0}+t)+f(x_{0}-t)-2s|\\
  &\leqslant|f(x_{0}+t)-f(x_{0}+0)|+|f(x_{0}-t)-f(x_{0}-0)|\\
  &<\frac{\varepsilon}{2}+\frac{\varepsilon}{2}=\varepsilon.
\end{align*}
\begin{align*}
  \left|\frac{1}{2n\pi}\int_{0}^{\delta}\varphi(t)\left[\frac{\sin\frac{nt}{2}}{\sin\frac{t}{2}}\right]^{2}\mathrm{d}t\right|
  &\leqslant\frac{1}{2n\pi}\int_{0}^{\delta}|\varphi(t)|\left[\frac{\sin\frac{nt}{2}}{\sin\frac{t}{2}}\right]^{2}\mathrm{d}t\\
  &<\frac{\varepsilon}{2n\pi}\int_{0}^{\delta}\left[\frac{\sin\frac{nt}{2}}{\sin\frac{t}{2}}\right]^{2}\mathrm{d}t
  <\frac{\varepsilon}{2n\pi}\int_{0}^{\pi}\left[\frac{\sin\frac{nt}{2}}{\sin\frac{t}{2}}\right]^{2}\mathrm{d}t=\frac{\varepsilon}{2},
\end{align*}
\begin{align*}
  \left|\frac{1}{2n\pi}\int_{\delta}^{\pi}\varphi(t)\left[\frac{\sin\frac{nt}{2}}{\sin\frac{t}{2}}\right]^{2}\mathrm{d}t\right|
  &\leqslant\frac{1}{2n\pi}\int_{\delta}^{\pi}|\varphi(t)|\left[\frac{\sin\frac{nt}{2}}{\sin\frac{t}{2}}\right]^{2}\mathrm{d}t\\
  &\leqslant\frac{1}{2n\pi\sin^{2}\frac{\delta}{2}}\int_{0}^{\pi}|\varphi(t)|\mathrm{d}t=\frac{A}{n},
\end{align*}
其中 $A=\dfrac{1}{2\pi\sin^{2}\frac{\delta}{2}}\displaystyle\int_{0}^{\pi}|\varphi(t)|\mathrm{d}t$ 为一个常数. 故当 $n>\dfrac{2A}{\varepsilon}$ 时,
\begin{align*}
  \left|\sigma_{n}(x_{0})-s\right|
  &=\left|\frac{1}{2n\pi}\int_{0}^{\delta}\varphi(t)\left[\frac{\sin\frac{nt}{2}}{\sin\frac{t}{2}}\right]^{2}\mathrm{d}t+\frac{1}{2n\pi}\int_{\delta}^{\pi}\varphi(t)\left[\frac{\sin\frac{nt}{2}}{\sin\frac{t}{2}}\right]^{2}\mathrm{d}t\right|\\
  &\leqslant\left|\frac{1}{2n\pi}\int_{0}^{\delta}\varphi(t)\left[\frac{\sin\frac{nt}{2}}{\sin\frac{t}{2}}\right]^{2}\mathrm{d}t\right|+\left|\frac{1}{2n\pi}\int_{\delta}^{\pi}\varphi(t)\left[\frac{\sin\frac{nt}{2}}{\sin\frac{t}{2}}\right]^{2}\mathrm{d}t\right|\\
  &<\frac{\varepsilon}{2}+\frac{A}{n}<\frac{\varepsilon}{2}+\frac{\varepsilon}{2}=\varepsilon.
\end{align*}
这就证明了 $f(x)$ 的 Fourier 级数在 $x_{0}$ 处的 Cesaro 和为
\[
  \lim_{n\to+\infty}\sigma_{n}=s=\frac{1}{2}\left[f(x_{0}+0)+f(x_{0}-0)\right].
\]
\end{proof}

\begin{corollary}\label{cor:16.4.1}
设 $f(x)$ 是周期为 $2\pi$ 且在 $[-\pi,\pi]$ 上可积或绝对可积的函数. 如果 $f(x)$ 在 $x_{0}$ 处有左、右极限, 且其 Fourier 级数在 $x_{0}$ 处收敛, 则必收敛于 $\dfrac{1}{2}[f(x_{0}+0)+f(x_{0}-0)]$.
\end{corollary}

\begin{proof}
设 $f(x)$ 的 Fourier 级数在 $x_{0}$ 处收敛于 $s$, 则其 Cesaro 和也必为 $s$. 根据定理~\ref{thm:16.4.2}, 其 Cesaro 和为 $\dfrac{1}{2}[f(x_{0}+0)+f(x_{0}-0)]$, 故得 $s=\dfrac{1}{2}[f(x_{0}+0)+f(x_{0}-0)]$.
\end{proof}

更进一步地, 有下面定理.

\begin{theorem}[Fej\'er 定理]\label{thm:16.4.3}
设 $f(x)$ 是周期为 $2\pi$ 的连续函数, 则它的 Fourier 级数在 Cesaro 意义下一致收敛于 $f(x)$.
\end{theorem}

\begin{proof}
设 $\varphi_{x}(t)=f(x+t)+f(x-t)-2f(x)$, 则
\[
  \sigma_{n}(x)-f(x)=\frac{1}{2n\pi}\int_{0}^{\pi}\varphi_{x}(t)\left[\frac{\sin\frac{nt}{2}}{\sin\frac{t}{2}}\right]^{2}\mathrm{d}t.
\]
由于 $f(x)$ 为整个数轴上的连续函数, 故在闭区间 $[-2\pi,2\pi]$ 上一致连续. 对于给定的 $\varepsilon>0$, $\exists\delta\in(0,\pi)$, 使得 $\forall t\in(0,\delta)$ 有
\[
  |f(x+t)-f(x)|<\frac{\varepsilon}{2},\quad|f(x-t)-f(x)|<\frac{\varepsilon}{2},\quad\forall x\in[-\pi,\pi].
\]
因而
\begin{align*}
  |\varphi_{x}(t)|&=|f(x+t)+f(x-t)-2f(x)|\\
  &\leqslant|f(x+t)-f(x)|+|f(x-t)-f(x)|\\
  &<\frac{\varepsilon}{2}+\frac{\varepsilon}{2}=\varepsilon,\quad\forall x\in[-\pi,\pi].
\end{align*}
于是, 当 $n>\dfrac{4M}{\varepsilon\sin^{2}\frac{\delta}{2}}$ 时(其中 $M=\max\limits_{x\in[-\pi,\pi]}|f(x)|$), 有
\begin{align*}
  |\sigma_{n}(x)-f(x)|
  &=\left|\frac{1}{2n\pi}\int_{0}^{\delta}\varphi_{x}(t)\left[\frac{\sin\frac{nt}{2}}{\sin\frac{t}{2}}\right]^{2}\mathrm{d}t+\frac{1}{2n\pi}\int_{\delta}^{\pi}\varphi_{x}(t)\left[\frac{\sin\frac{nt}{2}}{\sin\frac{t}{2}}\right]^{2}\mathrm{d}t\right|\\
  &\leqslant\frac{1}{2n\pi}\int_{0}^{\delta}|\varphi_{x}(t)|\left[\frac{\sin\frac{nt}{2}}{\sin\frac{t}{2}}\right]^{2}\mathrm{d}t+\frac{1}{2n\pi}\int_{\delta}^{\pi}|\varphi_{x}(t)|\left[\frac{\sin\frac{nt}{2}}{\sin\frac{t}{2}}\right]^{2}\mathrm{d}t\\
  &<\frac{\varepsilon}{2n\pi}\int_{0}^{\pi}\left[\frac{\sin\frac{nt}{2}}{\sin\frac{t}{2}}\right]^{2}\mathrm{d}t+\frac{2M}{n\sin\frac{\delta}{2}}<\frac{\varepsilon}{2}+\frac{\varepsilon}{2}=\varepsilon,
\end{align*}
因此
\[
  \sigma_{n}(x)\Rightarrow f(x),\quad x\in[-\pi,\pi],\ n\to+\infty.
\]
再由周期性, 有 $\sigma_{n}(x)\Rightarrow f(x)$, $x\in(-\infty,+\infty)$, $n\to+\infty$, 即 $f(x)$ 的 Fourier 级数在 Cesaro 意义下一致收敛于 $f(x)$.
\end{proof}

作为 Fej\'er 定理~\ref{thm:16.4.3} 与 Cesaro 意义下求和的应用, 我们来证明下面的 Weierstrass 定理.

\begin{theorem}[Weierstrass]\label{thm:16.4.4}
设 $f(x)$ 在 $[-\pi,\pi]$ 上连续, 且 $f(-\pi)=f(\pi)$, 则 $f(x)$ 必能用三角多项式一致逼近.
\end{theorem}

\textbf{证法 1}\quad 根据假定, 我们能将 $f(x)$ 延拓为整个数轴上的以 $2\pi$ 为周期的连续函数. 于是, 应用 Fej\'er 定理, $f(x)$ 能在 $(-\infty,+\infty)$ 中用序列 $\{\sigma_{n}(x)\}$ 一致逼近. 因为 $f(x)$ 的 Fourier 级数的 $k$ 次部分和 $S_{k}(x)$ 是一个 $k$ 次三角多项式, 因而
\[
  \sigma_{n}(x)=\frac{1}{n}\left[S_{0}(x)+S_{1}(x)+\cdots+S_{n-1}(x)\right]
\]
是一个 $n-1$ 次三角多项式, 它就是一个一致逼近 $f(x)$ 的三角多项式序列.

\textbf{证法 2}\quad 参阅定理~\ref{thm:14.3.8}.

% ----------------------------------------------------------------------
% 练习题 16.4
% ----------------------------------------------------------------------
\begin{exblock}{练习题 16.4}
\textbf{1.} 求下列级数的 Cesaro 和:
\begin{enumerate}
\item $1+0-1+1+0-1+\cdots$;
\item $\dfrac{1}{2}+\cos x+\cos 2x+\cdots+\cos nx+\cdots,\quad 0<x<2\pi$;
\item $\sin x+\sin 2x+\sin 3x+\cdots,\quad 0<x<2\pi.$
\end{enumerate}

\textbf{2.} 利用 Fej\'er 定理~\ref{thm:16.4.2} 证明: $[0,\pi]$ 上的连续函数可用余弦多项式一致逼近.

\textbf{3.} 证明: 级数 $\displaystyle\sum_{n=0}^{\infty}a_{n}$ 可以 Cesaro 求和的必要条件是 $a_{n}=o(n)$, $n\to+\infty$ (参阅思考题 12.1 题 5(3)).
\end{exblock}

% ----------------------------------------------------------------------
% 思考题 16.4
% ----------------------------------------------------------------------
\begin{exblock}{思考题 16.4}
\textbf{1.} 设由无穷级数 $\displaystyle\sum_{n=0}^{\infty}a_{n}$ 产生的幂级数 $\displaystyle\sum_{n=0}^{\infty}a_{n}x^{n}$ 的收敛半径为 $1$. 如果
\[
  \lim_{x\to1^{-}}\sum_{n=0}^{\infty}a_{n}x^{n}=s,
\]
则称级数 $\displaystyle\sum_{n=0}^{\infty}a_{n}$ 在 \textbf{Abel 意义下}收敛于 $s$, $s$ 称为级数 $\displaystyle\sum_{n=0}^{\infty}a_{n}$ 的 \textbf{Abel 和}, 记为 $\displaystyle\sum_{n=0}^{\infty}a_{n}=s(\mathrm{A})$. 此时, 称 $\displaystyle\sum_{n=0}^{\infty}a_{n}$ \textbf{可以 Abel 求和}.

\textbf{证明:} (1) 如果级数 $\displaystyle\sum_{n=0}^{\infty}a_{n}$ 在通常意义下收敛于 $s$, 则必有
\[
  \sum_{n=0}^{\infty}a_{n}=s(\mathrm{A}).
\]
(2) $\displaystyle\sum_{n=0}^{\infty}(-1)^{n}=\frac{1}{2}(\mathrm{A})$. 注意: $\displaystyle\sum_{n=0}^{\infty}(-1)^{n}=\frac{1}{2}(\mathrm{C})$, 但通常意义下它不收敛.

\textbf{2.} 证明: 如果 $\displaystyle\sum_{n=0}^{\infty}a_{n}=s(\mathrm{C})$, 则必有 $\displaystyle\sum_{n=0}^{\infty}a_{n}=s(\mathrm{A})$.

\textbf{提示:} 按下列步骤论证

(1) $\displaystyle\sum_{n=0}^{\infty}a_{n}x^{n}$ 在 $(-1,1)$ 中绝对收敛;

(2) 记 $f(x)=\displaystyle\sum_{n=0}^{\infty}a_{n}x^{n}$, 则
\[
  f(x)=(1-x)^{2}\sum_{n=0}^{\infty}(n+1)\sigma_{n+1}x^{n+1},\quad x\in(-1,1);
\]
(3) 从
\[
  f(x)-s=(1-x)^{2}\sum_{n=0}^{\infty}(n+1)(\sigma_{n+1}-s)x^{n}
\]
推得
\[
  \lim_{x\to1^{-}}\sum_{n=0}^{\infty}a_{n}x^{n}=s.
\]

\textbf{3.} 证明: 级数 $\displaystyle\sum_{n=0}^{\infty}(-1)^{n}(n+1)=\frac{1}{4}(\mathrm{A})$. 但它不能用 Cesaro 求和. 由此及题 1(2) 知, 通常求和最强, Cesaro 求和次强, Abel 求和最弱.

\textbf{4.} 证明: $\displaystyle\sum_{n=2}^{\infty}(-1)^{n}\ln n=\frac{1}{2}\ln\frac{\pi}{2}(\mathrm{C})$.
\end{exblock}
